\subsection{Introduction}

The success of any engineering project hinges on a clear, comprehensive understanding of the requirements that guide design, implementation, and evaluation. This chapter formalizes the requirement specification for the smart healthcare robot project, translating stakeholder needs and technical constraints into measurable, verifiable criteria. The requirements are systematically categorized into customer requirements, engineering requirements, and engineering constraints, ensuring alignment between user expectations and technical feasibility.

Customer requirements capture the high-level needs and expectations of primary stakeholders--hospital administrators, clinical staff, patients, and support personnel. These requirements reflect the functional and operational capabilities that directly impact system value and user satisfaction. Engineering requirements translate customer needs into quantifiable, testable technical specifications that guide system architecture and component selection. Engineering constraints define the boundaries within which the system must operate, including budget limitations, timeline restrictions, regulatory compliance, and environmental conditions specific to hospital deployment in Oman.

This chapter also presents a structured analysis of design alternatives, trade-offs, and selection criteria through comparative tables with metrics and scores. These analyses document the rationale for key architectural decisions, including mobile platform selection, sensor configuration, computing architecture, and AI integration strategy. The systematic evaluation ensures that design choices are transparent, justified, and optimized for the project's objectives and constraints.

By establishing a rigorous requirement framework, this chapter provides the foundation for subsequent design, implementation, and validation activities, ensuring that the final system meets stakeholder needs while remaining technically feasible and economically viable.

\subsection{Customer Requirements}

Customer requirements represent the high-level needs and expectations of the robot's primary users and stakeholders in hospital environments. These requirements are derived from literature review findings (Chapter 2), observations of existing hospital robotics deployments, and anticipated use cases in Omani healthcare facilities.

\begin{itemize}
    \item \textbf{CR1: Autonomous Navigation and Mobility.} The robot must autonomously navigate hospital corridors, patient rooms, and common areas without human intervention, avoiding static and dynamic obstacles. 
    
    \item \textbf{CR2: Remote Telepresence Capabilities.} The robot must enable high-quality remote video and audio communication between healthcare providers and patients, supporting remote consultations, rounds, and family visits. Telepresence has been validated as effective for remote diagnosis, specialist consultations, and reducing travel costs \cite{garingo2016}. The robot must provide a mobile "presence" for clinicians who are not physically at the patient's bedside.
    
    \item \textbf{CR3: Patient Engagement and Social Interaction.} The robot must engage with patients through natural language conversation, providing information, companionship, and emotional support, particularly for pediatric and elderly populations.
    
    \item \textbf{CR4: Infection Control and Hygiene.} The robot must be designed with materials and surfaces that are easy to clean and disinfect, minimizing cross-contamination risks. Hospital-acquired infections (HAIs) are a critical concern. Robots must be compatible with standard hospital cleaning protocols and potentially integrate UV-C disinfection capabilities \cite{intel2024, sciencedirect2021}.
    
    \item \textbf{CR5: User-Friendly Interface.} The robot must provide intuitive interfaces (touchscreen, mobile app, voice control) accessible to users with varying technical expertise. Ease of use directly impacts adoption rates and operational efficiency. Interfaces must accommodate diverse user groups including elderly patients, clinicians with minimal robotics experience, and administrative staff.
    
    \item \textbf{CR6: Continuous Operation and Autonomous Charging.} The robot must operate throughout the day with minimal downtime, autonomously returning to a charging dock when battery levels are low.
\end{itemize}

\subsection{Engineering Requirements}

Engineering requirements translate customer needs into quantifiable, testable technical specifications that directly inform system design and validation. Each engineering requirement is mapped to one or more customer requirements, ensuring traceability and alignment. The requirements are organized into five categories, presented in tabular format for clarity and traceability.

\subsubsection{Navigation and Mobility Requirements}

\begin{table}[H]
\centering
\small
\begin{tabularx}{\linewidth}{|l|>{\raggedright\arraybackslash}X|l|>{\raggedright\arraybackslash}X|>{\raggedright\arraybackslash}X|}
\hline
\textbf{Req. ID} & \textbf{Specification} & \textbf{Mapped to CR} & \textbf{Standards} & \textbf{Verification Method} \\
\hline
ER1.1 & Position accuracy of $\pm$20~cm in mapped indoor environments & CR1 & ISO 13482:2014 \cite{iso13482:2014} (Safety requirements for personal care robots - mobile servant robots) & Experimental validation using ground-truth positioning system \\
\hline
ER1.2 & Maximum linear velocity of 0.5--1.0~m/s in hospital corridors & CR1 & ISO 13482:2014 \cite{iso13482:2014} (Speed limits for safe human-robot interaction) & Speed measurement tests in controlled environments and hospital corridors \\
\hline
ER1.3 & Detect and avoid static and dynamic obstacles with minimum detection range of 3~m and reaction time $<$1~s & CR1 & ISO 13855 \cite{iso13855} (Safety of machinery - Positioning of safeguards) & Obstacle avoidance tests with various object types and approach angles \\
\hline
ER1.4 & Minimum payload of 5~kg to accommodate onboard computers, sensors, and delivery items & CR2 & --- & Load testing with representative payload configurations \\
\hline
\end{tabularx}
\caption{Navigation and mobility engineering requirements.}
\label{tab:er_navigation}
\end{table}

\subsubsection{Telepresence and Communication Requirements}

\begin{table}[H]
\centering
\small
\begin{tabularx}{\linewidth}{|l|>{\raggedright\arraybackslash}X|l|>{\raggedright\arraybackslash}X|>{\raggedright\arraybackslash}X|}
\hline
\textbf{Req. ID} & \textbf{Specification} & \textbf{Mapped to CR} & \textbf{Standards} & \textbf{Verification Method} \\
\hline
ER2.1 & Minimum 720p (1280$\times$720) video resolution at 30~fps & CR2 & ITU-T H.264 / H.265 \cite{itut_h264_h265} (Advanced Video Coding) & Video quality assessment tests under typical hospital lighting and network conditions \\
\hline
ER2.2 & Bidirectional audio communication with noise cancellation and echo reduction, maintaining intelligibility in hospital environments (ambient noise up to 70~dB) & CR2 & --- & Audio quality tests (speech intelligibility index, background noise rejection) \\
\hline
ER2.3 & End-to-end latency $<$200~ms over hospital Wi-Fi networks & CR2 & IEEE 802.11 \cite{ieee802.11} (Wi-Fi Standards) & Network performance testing and latency measurements \\
\hline
\end{tabularx}
\caption{Telepresence and communication engineering requirements.}
\label{tab:er_telepresence}
\end{table}

\subsubsection{Social Interaction and AI Requirements}

\begin{table}[H]
\centering
\small
\begin{tabularx}{\linewidth}{|l|>{\raggedright\arraybackslash}X|l|>{\raggedright\arraybackslash}X|>{\raggedright\arraybackslash}X|}
\hline
\textbf{Req. ID} & \textbf{Specification} & \textbf{Mapped to CR} & \textbf{Standards} & \textbf{Verification Method} \\
\hline
ER3.1 & Accurately interpret user intent in English and Arabic with $>$90\% accuracy for common healthcare queries & CR3 & --- & User testing with standardized query sets and intent classification metrics \\
\hline
ER3.2 & Generate conversational responses within 3~seconds of user input & CR3 & --- & Response latency measurements under typical operating conditions \\
\hline
ER3.3 & Speech recognition in hospital environments with word error rate (WER) $<$10\% & CR5 & --- & Speech recognition testing with diverse speakers and noise levels \\
\hline
\end{tabularx}
\caption{Social interaction and AI engineering requirements.}
\label{tab:er_ai}
\end{table}

\subsubsection{Power and Endurance Requirements}

\begin{table}[H]
\centering
\small
\begin{tabularx}{\linewidth}{|l|>{\raggedright\arraybackslash}X|l|>{\raggedright\arraybackslash}X|>{\raggedright\arraybackslash}X|}
\hline
\textbf{Req. ID} & \textbf{Specification} & \textbf{Mapped to CR} & \textbf{Standards} & \textbf{Verification Method} \\
\hline
ER4.1 & Continuous operation for minimum of 4--6~hours under typical usage (mixed navigation, telepresence, and idle states) & CR6 & --- & Endurance testing with simulated clinical workflows \\
\hline
ER4.2 & Autonomously locate and dock with charging station when battery level reaches 20\%, with 95\% success rate & CR6 & --- & Repeated docking trials under various initial positions and orientations \\
\hline
\end{tabularx}
\caption{Power and endurance engineering requirements.}
\label{tab:er_power}
\end{table}

\subsubsection{Safety and Reliability Requirements}

\begin{table}[H]
\centering
\small
\begin{tabularx}{\linewidth}{|l|>{\raggedright\arraybackslash}X|l|>{\raggedright\arraybackslash}X|>{\raggedright\arraybackslash}X|}
\hline
\textbf{Req. ID} & \textbf{Specification} & \textbf{Mapped to CR} & \textbf{Standards} & \textbf{Verification Method} \\
\hline
ER5.1 & Physical emergency stop button(s) and remote e-stop capability via software interface & CR1 & ISO 13850 \cite{iso13850} (Safety of machinery - Emergency stop function) & Functional testing of emergency stop mechanisms \\
\hline
ER5.2 & Detect and avoid stairs, ramps $>$15° incline, and drop-offs $>$2~cm & CR1 & --- & Edge detection tests on stairs and elevated platforms \\
\hline
ER5.3 & Achieve $>$95\% uptime during operational hours, excluding scheduled maintenance & CR6 & --- & Long-term reliability testing and failure logging \\
\hline
\end{tabularx}
\caption{Safety and reliability engineering requirements.}
\label{tab:er_safety}
\end{table}

\subsubsection{Comparative Performance Metrics}

This section presents quantitative performance data and trade-off metrics for critical system components. The data, derived from technical specifications and benchmark studies, provides the empirical basis for architectural decisions. Comprehensive mobile platform comparison metrics, including navigation capabilities, payload specifications, and integration features, are detailed in the literature review (Chapter 2, Section 2.3.2).

\paragraph{Navigation Sensor Performance}
Table~\ref{tab:sensor_metrics} compares the performance of different sensor modalities in typical hospital environments.

\begin{table}[H]
\centering
\scriptsize
\begin{tabularx}{\linewidth}{|>{\raggedright\arraybackslash}X|>{\raggedright\arraybackslash}X|c|c|c|c|c|}
\hline
\textbf{Sensor Type} & \textbf{Scenario} & \textbf{Range (m)} & \textbf{Accuracy (cm)} & \textbf{FOV (deg)} & \textbf{Power (W)} & \textbf{Cost (\$)} \\
\hline
2D LiDAR (RPLiDAR) \cite{rplidar2024} & Corridor (Static) & 12.0 & $\pm$1.0 & 360 & 1.5 & 100 \\
\hline
2D LiDAR (RPLiDAR) \cite{rplidar2024} & Waiting Room (Dynamic) & 10.0 & $\pm$2.5 & 360 & 1.5 & 100 \\
\hline
RGB-D (RealSense) \cite{realsense2024} & Corridor (Static) & 4.0 & $\pm$2.0 & 87$\times$58 & 3.5 & 250 \\
\hline
RGB-D (RealSense) \cite{realsense2024} & Waiting Room (Dynamic) & 3.5 & $\pm$4.0 & 87$\times$58 & 3.5 & 250 \\
\hline
3D LiDAR (Velodyne) \cite{velodyne2024} & Corridor (Static) & 100.0 & $\pm$2.0 & 360$\times$30 & 8.0 & 4000 \\
\hline
\end{tabularx}
\caption{Performance metrics of navigation sensors in static and dynamic scenarios.}
\label{tab:sensor_metrics}
\end{table}

\paragraph{Computing Platform Benchmarks}
Table~\ref{tab:compute_metrics} details the computational capabilities and power efficiency of candidate onboard computers.

\begin{table}[H]
\centering
\footnotesize
\begin{tabular}{|l|c|c|c|}
\hline
\textbf{Platform} & \textbf{CPU Perf.} & \textbf{Power (W)} & \textbf{Cost (\$)} \\
\hline
Raspberry Pi 5 \cite{raspberrypi2024} & High & 5--8 & 80 \\
Jetson Nano \cite{jetsonnano2024} & Medium & 5--10 & 150 \\
Jetson Orin Nano \cite{jetsonorin2024} & Very High & 7--15 & 499 \\
\hline
\end{tabular}
\caption{Computational performance and power consumption benchmarks for embedded platforms.}
\label{tab:compute_metrics}
\end{table}

\paragraph{Conversational AI Model Selection}
Detailed performance and cost metrics for Large Language Models (LLMs), including context handling capabilities and operational costs for cloud-based and local deployment strategies, are presented in the literature review (Chapter 2, Section 2.3.3). These metrics inform the AI architecture selection for patient engagement and conversational interaction requirements (ER3.1--ER3.3).

\subsection{Engineering Constraints}

Engineering constraints define the boundaries within which the system must be designed and operated. These constraints reflect real-world limitations in budget, timeline, technology maturity, regulatory environment, and deployment context specific to Oman.

\begin{itemize}
    \item \textbf{EC1: Budget Constraints.} Total system cost must not exceed \$2,500 USD, including mobile platform, sensors, computing hardware, and accessories. As an educational research project with limited funding, the system must demonstrate feasibility and functionality within academic budget constraints while remaining representative of scalable commercial solutions. This limits selection to budget and mid-tier commercial platforms; excludes high-end industrial robots; requires prioritization of essential features.
    
    \item \textbf{EC2: Development Timeline.} System design, implementation, and validation must be completed within one academic year (9--10 months). Final year project timeline constraints necessitate realistic scope and phased development approach. This requires prioritization of core functionalities (navigation, telepresence, basic AI); potential deferral of advanced features (UV-C disinfection, complex HIS integration) to future work.
    
    \item \textbf{EC3: Hospital Environment Compatibility.} The robot must operate safely in hospital environments with corridor widths $\geq$1.5~m, elevator access, and typical hospital flooring (tile, linoleum). Omani hospital infrastructure follows international standards with standard corridor dimensions and flooring materials. Robot base width limited to $<$60~cm; wheel system must handle smooth indoor surfaces; turning radius must accommodate corridor navigation.
    
    \item \textbf{EC4: Network Infrastructure.} The robot must operate on standard hospital Wi-Fi networks (802.11n/ac) without requiring dedicated network infrastructure. Deployment feasibility requires compatibility with existing IT infrastructure; hospitals cannot always support custom networking solutions. System must tolerate variable network quality, latency, and temporary disconnections; requires local processing capabilities for critical functions.
    
    \item \textbf{EC5: Regulatory and Privacy Compliance.} The system must comply with patient privacy requirements and avoid collection or storage of protected health information (PHI) without proper authorization. Healthcare robotics must respect patient privacy and data security regulations (similar to HIPAA/GDPR principles). No persistent storage of patient conversations or video recordings; encrypted communication channels; clear data handling policies.
    
    \item \textbf{EC6: User Expertise Level.} The system must be operable by users with minimal technical training (hospital staff, patients) without requiring robotics expertise. Clinical adoption depends on ease of use and minimal training burden. Intuitive UI design; pre-configured waypoints; automated error recovery; comprehensive user documentation.
    
    \item \textbf{EC7: Environmental Operating Conditions.} The robot must operate reliably in indoor environments with temperature range 18--28°C, humidity 30--60\%, and typical hospital lighting conditions. Hospital climate control maintains stable environmental conditions; sensors and electronics must function within this range. Standard commercial-grade components acceptable; no specialized environmental hardening required.
    
    \item \textbf{EC8: Maintenance and Spare Parts.} The robot must utilize commercially available, globally sourced components with accessible supply chains for replacement parts. Long-term operational viability requires availability of replacement components in Oman or via international suppliers. Preference for mainstream platforms (ROS 2, Raspberry Pi, common sensors) over proprietary or region-specific components.
\end{itemize}

\subsection{Analysis and Discussion}

This section examines the comparative performance metrics to identify key trade-offs and demonstrate standards compliance within the constraints of hospital deployment in Oman. The analysis focuses on component selection criteria that balance technical requirements with practical deployment considerations specific to the Omani healthcare context.

\subsubsection{Performance Metrics and Trade-off Analysis}

Table~\ref{tab:sensor_metrics} reveals fundamental trade-offs in sensor selection. Range capabilities vary from 3.5m to 100m with corresponding cost differences of 40$\times$ (\$100 to \$4000). Accuracy degrades significantly in dynamic hospital environments (2--2.5$\times$ degradation), affecting position accuracy requirement ER1.1. Power consumption ranges from 1.5W to 8W, directly impacting endurance requirement ER4.1. Field of view characteristics--360° planar, 87°$\times$58° 3D, and 360°$\times$30° volumetric--address different aspects of obstacle detection requirement ER1.3.

Table~\ref{tab:compute_metrics} demonstrates computational trade-offs where platform costs range from \$80 to \$499, consuming 3--20\% of total budget constraint EC1. Power envelopes remain similar (5--15W) across platforms, making computational capability rather than endurance the primary selection factor for supporting navigation, AI processing (ER3.1--ER3.3), and telepresence (ER2.1--ER2.3) requirements.

\subsubsection{Standards Compliance and Deployment Context}

Navigation and safety components satisfy ISO 13482:2014 \cite{iso13482:2014} requirements (position accuracy ER1.1, speed limits ER1.2, obstacle detection ER1.3), ISO 13850 \cite{iso13850} emergency stop guidelines, and ISO 13855 \cite{iso13855} safeguarding standards. Communication systems implement ITU-T H.264/H.265 \cite{itut_h264_h265} video coding and IEEE 802.11 \cite{ieee802.11} wireless protocols for interoperability with Omani hospital IT infrastructure.

Deployment constraints EC3--EC8 address Omani hospital realities: standard corridor dimensions ($\geq$1.5m) and flooring materials, existing Wi-Fi infrastructure (802.11n/ac), climate-controlled environments (18--28°C), component availability through international supply chains, privacy compliance aligned with HIPAA/GDPR principles, and budget/timeline limitations (\$2,500, 9--10 months). These constraints guide component selection to ensure technical feasibility and operational viability in the Omani healthcare context.

Detailed comparative analysis of mobile platforms, computing architectures, and conversational AI models is presented in the literature review (Chapter 2, Sections 2.3.2--2.3.3).

\subsection{Conclusion}

This chapter established a comprehensive requirement specification for the smart healthcare robot project, systematically translating stakeholder expectations into measurable, verifiable technical criteria. The requirements are organized into customer requirements, engineering requirements, and engineering constraints, defining the system's functional capabilities, performance targets, and operational boundaries.

Engineering requirements reference relevant international standards including ISO 13482:2014 (personal care robots), ISO 13855 and ISO 13850 (machinery safety), ITU-T H.264/H.265 (video coding), and IEEE 802.11 (wireless communications), establishing the regulatory framework for system validation. Engineering constraints reflect budget limitations, timeline restrictions, and regulatory compliance needs specific to Omani healthcare deployment.

The comparative performance analysis in Section 3.3.6 provided empirical data on component characteristics, while Section 3.5 examined trade-offs in the context of standards compliance and Omani deployment requirements. The requirement specification establishes the foundation for system design decisions detailed in subsequent chapters, with traceability ensuring design choices can be validated against stakeholder needs and technical constraints.
